\documentclass[12pt,onecolumn,letterpaper]{article}

\usepackage{geometry}
\geometry{letterpaper, left=0.75in, right=0.75in, bottom=0.75in, top=0.95in}

\usepackage{fontspec}
\usepackage{microtype}
\usepackage{setspace}
\usepackage{array}
\usepackage{booktabs}
\usepackage{tabularx}
\usepackage{hyperref}
\usepackage{amsmath}
\usepackage{listings}
\usepackage[table]{xcolor}
\usepackage{caption}
\usepackage{enumitem}
\usepackage{appendix}
\usepackage{titlesec}
\usepackage{fancyhdr}
\usepackage{indentfirst}

\setmainfont{Cambria}
\setmonofont{Liberation Mono}[Scale=MatchLowercase]
\setlength{\parindent}{0pt}
\setlength{\parskip}{0.45em}
\setstretch{1.08}
\setlength{\headheight}{15pt}
\setlength{\emergencystretch}{2em}
\widowpenalty=10000
\clubpenalty=10000
\displaywidowpenalty=10000

\definecolor{MEBlue}{HTML}{1E3A5F}
\definecolor{MEGray}{HTML}{E3E8ED}
\definecolor{MERule}{HTML}{D6DADF}

\newcommand{\sectionbreak}{\clearpage}

\titleformat{\section}{\Large\bfseries}{\thesection}{0.65em}{}
\titleformat{\subsection}{\large\bfseries}{\thesubsection}{0.6em}{}
\titleformat{\subsubsection}{\normalsize\bfseries}{\thesubsubsection}{0.5em}{}
\titlespacing*{\section}{0pt}{1.55em}{0.7em}
\titlespacing*{\subsection}{0pt}{1.15em}{0.45em}
\titlespacing*{\subsubsection}{0pt}{0.9em}{0.35em}

\pagestyle{fancy}
\fancyhf{}
\fancyhead[L]{\small\itshape MinEnglish Language Specification}
\fancyhead[R]{\small\thepage}
\renewcommand{\headrulewidth}{0.35pt}
\renewcommand{\footrulewidth}{0pt}

\urlstyle{same}
\hypersetup{
    colorlinks=true,
    linkcolor=MEBlue,
    filecolor=MEBlue,
    urlcolor=MEBlue,
    pdftitle={MinEnglish Language Specification V1.0.0},
    pdfauthor={Dan Micsa, PhD},
}

\lstset{
    basicstyle=\ttfamily\small,
    breaklines=true,
    frame=single,
    backgroundcolor=\color{gray!10},
    keywordstyle=\bfseries,
    commentstyle=\itshape
}

\captionsetup[table]{font=small,labelfont=bf,labelsep=period,skip=6pt,justification=raggedright,singlelinecheck=false}

\newcommand{\me}[1]{\texttt{#1}}
\newcommand{\mecode}[1]{\texttt{\detokenize{#1}}}
\newcolumntype{Y}{>{\raggedright\arraybackslash}X}
\newcommand{\MEtablelayout}{%
\footnotesize
\setlength{\tabcolsep}{6pt}%
\renewcommand{\arraystretch}{1.12}%
\arrayrulecolor{MERule}%
}
\newcommand{\corpusExample}[8]{%
\subsubsection*{Ex.~#1}
\begin{center}
\MEtablelayout
\begin{tabularx}{\linewidth}{@{}>{\bfseries}p{2.35cm}Y>{\raggedleft\arraybackslash}p{1.1cm}@{}}
\rowcolor{MEGray}
 & \textbf{Text} & \textbf{Chars} \\
English & \mecode{#2} & #5 \\
MinEnglish & \mecode{#3} & #6 \\
Back & \mecode{#4} & \\
\bottomrule
\end{tabularx}
\end{center}
\noindent\textbf{Compression:} #7\% $\Delta$\\
\textbf{Observation:} #8\par\vspace{0.85em}
}

\title{\Huge\textbf{MinEnglish Language Specification}\\[0.2em]
\LARGE\textit{V1.0.0 --- Stable Release}\\[0.3em]
\large A Zero-Redundancy Phonetic Communication Protocol}
\author{\textbf{Dan Micsa, PhD}\\
\small MIT License (see Appendix)}
\date{March 2026}

\begin{document}

\maketitle
\thispagestyle{empty}

%--- AI ATTRIBUTION
\begin{quote}
\small\textit{AI Attribution: This specification was developed with extensive
assistance from large language models. \textbf{Claude Opus 4.6 Thinking}
(Anthropic), \textbf{Gemini 3.1 Pro} (Google DeepMind), and \textbf{GPT-5.4}
(OpenAI) contributed to the drafting, formalization, revision, and example
generation documented herein. The intellectual framework, design decisions,
and authorship remain those of the human author.}
\end{quote}

\begin{abstract}
MinEnglish is a formal communication protocol engineered to maximize Shannon
information density$^{[1]}$ while minimizing cognitive load and parser
ambiguity. By systematically eliminating morphological irregularities and
orthographic redundancies inherent in natural languages, MinEnglish achieves
consistent computational predictability. This specification (V1.0.0) defines
the phonological, morphological, and syntactic architecture targeting the
V1.0.0 stable release, documents a 35-sentence comparative corpus
demonstrating an average 28.2\% character reduction over standard American
English, and addresses the primary theoretical limitation of the system:
biological adoption friction.
\end{abstract}

\tableofcontents
\vspace{1em}

%---
\section{Introduction: Design Principles}

MinEnglish is governed by five absolute constraints, each eliminating a
specific failure mode of natural language:

\begin{itemize}[leftmargin=*]
    \item \textbf{Phonetic Isomorphism:} 1-to-1 grapheme-phoneme correspondence eliminates orthographic memorization overhead and provides predictable tokenization for machine parsers.
    \item \textbf{Morphological Parsimony:} Pluralization, conjugation, and declension are replaced by explicit numerical and temporal prefixes, removing morphological boundary ambiguity.
    \item \textbf{ASCII Strictness:} Operates exclusively within the standard US keyboard namespace without diacritics. Letters \textit{q}, \textit{w}, and \textit{y} are omitted to balance acoustic distinctness with key-space economy.
    \item \textbf{Operator Integration:} Symbolic tokens (\me{*}, \me{\~{}}, \me{>}, \me{<}, \me{:}, \me{\&}, \me{|}) function as first-class syntactic elements, compressing multi-word constructs into single-byte representations.
    \item \textbf{Absolute Regularity:} A zero-exception framework guarantees programmatic rules apply universally, ensuring high predictability for both human acquisition and NLP models$^{[2]}$.
\end{itemize}

\textbf{Versioning:} MinEnglish follows Semantic Versioning (SemVer):
\texttt{MAJOR.MINOR.PATCH}. Breaking rule changes increment MAJOR.
\texttt{0.x.y} indicates pre-stable development; \texttt{1.0.0} marks the
first frozen specification.

%---
\section{Phonology and Orthography}

\subsection{Vowels}
The vowel space is constrained to five cardinal points for acoustic
discernibility across diverse linguistic backgrounds.

\begin{table}[h!]
\centering
\MEtablelayout
\begin{tabularx}{\linewidth}{lllX}
\rowcolor{MEGray}
\textbf{Graph.} & \textbf{Sound} & \textbf{English} & \textbf{MinEnglish} \\
\textbf{a} & as in c\textit{a}t & cat & cat \\
\textbf{e} & as in b\textit{e}d & bed & bed \\
\textbf{i} & as in b\textit{i}t & bit & bit \\
\textbf{o} & as in h\textit{o}t & hot & hot \\
\textbf{u} & as in b\textit{u}t & but & but \\
\bottomrule
\end{tabularx}
\caption{Cardinal Vowel Mapping}
\end{table}

\subsection{Long Vowels and Diphthongs}
Diphthongs are encoded as contiguous grapheme pairs, preserving the 1-to-1
mapping without macron diacritics: \me{ei}~(day), \me{ii}~(see),
\me{ai}~(like), \me{ou}~(go), \me{uu}~(food), \me{au}~(how),
\me{oi}~(boy), \me{ar}~(car), \me{or}~(for), \me{er}~(her).

\subsection{Consonants}

\begin{table}[h!]
\centering
\MEtablelayout
\begin{tabularx}{\linewidth}{llX}
\rowcolor{MEGray}
\textbf{Graph.} & \textbf{Sound} & \textbf{Rule / Example} \\
\textbf{c}  & /k/ soft      & Unaspirated velar: cat, cum, bac \\
\textbf{k}  & /k/ hard      & Aspirated velar: king, kiip, kic \\
\textbf{ci} & /ch/          & Affricate (church): ciald, muci \\
\textbf{u}  & /w/           & Semivowel for \textit{w}: uoter \\
\textbf{i}  & /y/           & Semivowel for \textit{y}: ies, iir \\
\textbf{sh} & /sh/          & As in ship, cash \\
\textbf{zh} & /zh/          & As in vision: vizhun \\
\bottomrule
\end{tabularx}
\caption{Non-standard Consonant Mappings}
\end{table}

\textbf{Aspiration (\me{k} vs \me{c}):} Aspirated /k\textsuperscript{h}/
uses \me{k}; unaspirated uses \me{c}. Following sibilants, \me{c} is
mandated (e.g., \me{scan}, \me{scip}).

\subsection{Orthographic Normalization}
Historical artifacts are removed algorithmically: silent letters nullified,
\me{-tion}\,$\to$\,\me{-shun}, \me{ck} resolved by aspiration,
\me{x}\,$\to$\,\me{cs}, double consonants collapsed, soft \me{c}\,$\to$\,\me{s}.

%---
\section{Morphology and Syntax}

\textbf{Core Constraint:} Lexical items never undergo morphological inflection.
Grammatical relationships are managed exclusively through operators and
rigid syntactic ordering.

\subsection{Nouns and Quantitative Prefixes}
Morphological plurals are abolished. Quantity is mandated explicitly:

\begin{table}[h!]
\centering
\MEtablelayout
\begin{tabularx}{\linewidth}{lXl}
\rowcolor{MEGray}
\textbf{Prefix} & \textbf{Semantic Value} & \textbf{Example} \\
\me{1}     & Singular / Unitary  & \me{1cat} \\
\me{3}     & Exact enumeration   & \me{3cat} \\
\me{*}     & Universal / Generic & \me{*cat} (all cats) \\
\me{\~{}5} & Approximate         & \me{\~{}5cat} \\
\me{0}     & Null set            & \me{0cat} \\
\bottomrule
\end{tabularx}
\caption{Quantitative Prefix System}
\end{table}

\subsubsection{Token Optionality Rules (V1.0.0)}
Three token classes may be omitted when recoverable from context:
\begin{itemize}[leftmargin=*]
    \item \textbf{Rule A --- Implicit Singular:} Bare noun without prefix defaults to singular. \me{cat sit on mat} = \textit{A cat sits on a mat.}
    \item \textbf{Rule B --- Implicit First Person:} When no subject precedes the verb, parser defaults to \me{i}. \me{laic cat} = \textit{I like a cat.}
    \item \textbf{Rule C --- Implicit Present:} A bare verb with no temporal prefix is present continuous. \me{h laic cat} = \textit{He likes a cat.}
\end{itemize}

\subsection{Compound Word Operators (\me{-} and \me{\_})}

\textbf{Hyphen \me{-} --- directional ``of/for'' binder:}
Left-to-right ``X for/of Y'' reading.
\begin{itemize}[leftmargin=*,noitemsep]
    \item \me{haus-fuud} = house \textit{for} food (restaurant)
    \item \me{haus-sic} = house \textit{of} the sick (hospital)
    \item \me{farmasi} = pharmacy (canonical phonetic adaptation)
\end{itemize}

\textbf{Disambiguation:} When \me{-} precedes a digit or \me{0s}, it is a
temporal operator. When it binds two noun roots, it is a compound operator.

\textbf{Underscore \me{\_} --- literalization:} Forces strict compositional
reading. \me{haus\_fuud} = a house built from food (gingerbread house),
vs.\ \me{haus-fuud} = restaurant (cultural abstraction).

\textbf{Note:} The \me{/} character is reserved exclusively for fractions
and mathematical division (\me{1/2}, \me{3/4}). It is not a compound operator.

\subsection{Proper Noun Encapsulation}
Proper nouns retain canonical native orthography. Initial capitalization acts
as an escape character suspending local phonetic rules: \me{person Isabella},
\me{cuntri France}.

\subsection{Mathematical Notation}
MinEnglish adopts LaTeX inline notation as a first-class citizen. Tokens
enclosed in \texttt{\$...\$} are parsed as LaTeX and are exempt from phonetic
rules: \texttt{\$E = mc\^{}2\$}, \texttt{\$\textbackslash frac\{1\}\{2\}\$}.
Spoken form is conventional English pronunciation.

\subsection{Verbs and Temporal Mechanics}

\begin{table}[h!]
\centering
\MEtablelayout
\begin{tabularx}{\linewidth}{lX}
\rowcolor{MEGray}
\textbf{Operator} & \textbf{Temporal Value / Example} \\
\textit{(null)} & Present continuous: \me{iit} (eating) \\
\me{1}          & Present discrete: \me{i1iit} (I eat) \\
\me{*}          & Habitual/Universal: \me{*iit} (eats always) \\
\me{-0s}        & Immediate past: \me{-0s sii} (just saw) \\
\me{-1d}        & Past offset: \me{-1d iit} (ate yesterday) \\
\me{:5Y}        & Duration scalar: \me{:5Y studi} (studied for 5 yrs) \\
\me{1d}         & Future offset (\me{+} optional): \me{1d gou} (will go) \\
\me{2026-12-25} & Absolute timestamp \\
\bottomrule
\end{tabularx}
\caption{Temporal Prefix System (V1.0.0)}
\end{table}

\textbf{Direction Convention:} The \me{-} prefix marks past vectors
(required). Future offsets are positive by default --- a bare numeric offset
(\me{1d}, \me{2h}, \me{1Y}) always points forward. The \me{+} prefix is
accepted by parsers but considered redundant.

\textbf{Temporal Stacking:} Left-to-right chaining enables algebraic
timeline construction: \me{-0s -1d iit} = ``I had already eaten by
yesterday's timeline.''

\subsection{Symbolic Operator Inventory}

\begin{table}[h!]
\centering
\MEtablelayout
\begin{tabularx}{\linewidth}{llX}
\rowcolor{MEGray}
\textbf{Op.} & \textbf{Name} & \textbf{Role / Example} \\
\me{+}  & Plus      & Future offset (optional): \me{1d gou} \\
\me{-}  & Minus     & Past offset / compound binder: \me{-1d}, \me{haus-fuud} \\
\me{*}  & Star      & Universal quantifier: \me{*cat} \\
\me{\&} & Ampersand & Logical AND ($\equiv$ \me{an}): \me{cat \& dog} \\
\me{|}  & Pipe      & Logical OR ($\equiv$ \me{or}): \me{cat | dog} \\
\bottomrule
\end{tabularx}
\caption{Symbolic Operators (all single-byte ASCII)}
\end{table}

Note: \me{*i} = ``all of us'' (universal quantifier on pronoun), not
habitual verb. Parsers resolve by token class.

\subsection{Scalar Intensity}
Adjectives post-modify nouns. Intensity uses \me{>} and \me{<} combinators:
\me{>big} (very big), \me{>>big} (extremely big), \me{<hot} (slightly hot).

\subsection{Pronominal Indexing}
Single-character spatial indices: \me{i}~(I), \me{u}~(you), \me{h}~(he),
\me{s}~(she), \me{t}~(it). Plurality inherits the quantitative prefix system
(\me{5i} = we five, \me{*i} = all of us). Pronominal index agglutinates to
the temporal operator: \me{h1laic}, \me{s-1d gou}, \me{u:2h uorc}.

Possession via \me{'s} suffix: \me{i's} (my), \me{h's} (his), \me{s'} (her).

\subsection{Negation}
Invariant prefix \me{no} prepended to the verb: \me{no laic cat} (I don't
like the cat). The Focus Operator \me{!} targets scope: \me{no laic !3cat}
(not \textit{three} cats; I like four).

\subsection{Modality}
\me{can} (ability), \me{must} (obligation), \me{shud} (advice),
\me{\~{}can} (probability). Attach to temporal block: \me{u 1d must cum}
(you must come tomorrow), \me{i 1d \~{}can gou} (I might go tomorrow).

\subsection{Passive Voice}
Object fronted, verb followed by \me{bai}: \me{bol kic bai boi} (ball
kicked by boy).

\subsection{Connectors}
\me{an}/\me{\&}~(and), \me{or}/\me{|}~(or), \me{but}, \me{if}, \me{cos}
(because), \me{sou}~(so), \me{den}~(then). Deference: \me{\~{}u giv}
= ``could you please give.''

\subsection{Clause Anchors (\me{co} / \me{oc})}
When nesting 2+ clauses, \me{co} (clause open) and \me{oc} (clause close)
act as explicit stack managers, preventing center-embedding overload of
the human phonological loop$^{[3]}$:\\
\me{man co dat dog co dat ciald luv oc -1d bait oc run auei.}

\subsection{Echo-Tagging (Noisy Protocol)}
Per Shannon's Theorem$^{[1]}$, redundancy may be inflated to 100\% for
critical transmissions by echoing prefixes as full words: \me{3boi trii uoc nau}.

\subsection{Formal Syntax (EBNF)}

\begin{lstlisting}[language=SQL]
<sentence> ::= [<subj>] <verb_block> [<obj>]
<subj>     ::= <noun_phrase> | <pronoun>
<noun_phrase> ::= <qty> <root> [<adj>]
<verb_block>  ::= [<temporal>] [<modal>] <root>
<temporal>    ::= <offset> | <duration> | <timestamp>
<clause>      ::= "co" <sentence> "oc"
\end{lstlisting}

%---
\section{Base Lexicon}

\begin{table}[h!]
\centering
\MEtablelayout
\begin{tabularx}{\linewidth}{lXlX}
\rowcolor{MEGray}
\textbf{English} & \textbf{MinEnglish} & \textbf{English} & \textbf{MinEnglish} \\
be / is     & \me{bi}   & have       & \me{hav}   \\
do          & \me{du}   & go         & \me{gou}   \\
come        & \me{cum}  & make       & \me{meic}  \\
think       & \me{tinc} & say        & \me{sei}   \\
want        & \me{uant} & need       & \me{niid}  \\
see         & \me{sii}  & know       & \me{nou}   \\
like        & \me{laic} & give       & \me{giv}   \\
eat         & \me{iit}  & run        & \me{run}   \\
\bottomrule
\end{tabularx}
\caption{Core Fundamentals}
\end{table}

\subsection{Academic \& Computational Lexicon}

\begin{table}[h!]
\centering
\MEtablelayout
\begin{tabularx}{\linewidth}{llX}
\rowcolor{MEGray}
\textbf{Domain} & \textbf{English} & \textbf{MinEnglish (Roots)} \\
\textit{Science}    & Hypothesis  & \me{aidia-test} (idea-test) \\
                    & Algorithm   & \me{rul-mat} (rule-math) \\
                    & Evolution   & \me{ceinj-taim} (change-time) \\
                    & Gravity     & \me{pul-erd} (pull-earth) \\
\midrule
\textit{Computing}  & Database    & \me{daatabeis} (phonetic) \\
                    & Variable    & \me{ting-ceinj} (thing-change) \\
                    & Memory      & \me{haus-daata} (house-data) \\
                    & Network     & \me{net-lain} (net-line) \\
\midrule
\textit{Economics}  & Inflation   & \me{rap-mun} (rapid-money) \\
                    & Contract    & \me{rul-peper} (rule-paper) \\
                    & Tax         & \me{tacs} (phonetic) \\
\midrule
\textit{Medicine}   & Antibiotic  & \me{bac-med} (bacteria-medicine) \\
                    & Diagnosis   & \me{faind-sic} (find-sick) \\
                    & Hospital    & \me{haus-sic} (house-of-sick) \\
                    & Pharmacy    & \me{farmasi} (phonetic adapt.) \\
\bottomrule
\end{tabularx}
\caption{Academic Domain Lexicon (Representative Sample)}
\end{table}

%---
\section{Comparative Corpus (n=35)}

The following 35 examples demonstrate MinEnglish performance across diverse
linguistic domains. Back-translations reconstruct compositional logic.
Compression = character count reduction vs. standard English.

\subsection{Simple Statements \& Universal Truths}

\corpusExample{1}
{The three cats are sitting on the big red mat.}
{3cat sit on mat big red.}
{Three-cat sit on one-mat big red.}
{46}{24}{47.8}
{High compression comes from removing articles, auxiliary verbs, and redundant inflection.}

\corpusExample{2}
{Dogs are loyal animals that love their owners.}
{*dog bi *animal loial dat *luv *h's ouner.}
{Any-dog be any-animal loyal that any-love any-their owner.}
{46}{42}{8.7}
{Universal quantification keeps the meaning precise, but the sentence still depends on longer lexical material.}

\subsection{Past Events, Times, and Futures}

\corpusExample{3}
{Yesterday I bought two books and three apples at the store.}
{-1d bai 2buc an 3apel @ stor.}
{I one-day-ago buy two-book and three-apple at one-store.}
{59}{29}{50.8}
{Temporal prefixes compress tense and adverbial framing into a short operator block.}

\corpusExample{4}
{Will you eat dinner with us tomorrow evening?}
{u1d iit diner uid *i?}
{You in-one-day eat one-dinner with all-of-us?}
{45}{21}{53.3}
{The future operator and compact pronoun group remove a large auxiliary phrase.}

\corpusExample{5}
{About fifty people arrived at approximately three o'clock.}
{~50person ariv @ ~15:00.}
{Approximately-fifty-person arrive at approximately-15:00.}
{58}{24}{58.6}
{Approximation and clock notation replace long English paraphrases very efficiently.}

\corpusExample{6}
{The meeting was held on January fifteenth, twenty twenty-six, at two thirty in the afternoon.}
{miiting 2026-01-15 14:30 held.}
{One-meeting 2026-01-15 14:30 held.}
{93}{30}{67.7}
{Absolute timestamp notation produces the strongest compression in the corpus.}

\corpusExample{7}
{I have been studying English for five years but I still cannot speak fluently.}
{:5Y studi inglish but i no can spiic fluuent.}
{For-five-years study English but I not can speak fluent.}
{78}{45}{42.3}
{Duration marking removes a long perfect-progressive construction with one compact prefix.}

\subsection{Complex Modifiers \& Voice}

\corpusExample{8}
{We are going to the extremely new restaurant tomorrow at seven in the evening.}
{*i1d gou tu restaurant >>niu @ 19:00.}
{All-of-us in-one-day go to one-restaurant extremely-new at 19:00.}
{78}{37}{52.6}
{Intensity markers and time notation shorten descriptive detail without changing meaning.}

\corpusExample{9}
{Look out! The glass window was just broken by the angry man.}
{uac!! uindou glas -0s breic bai man angri.}
{Watch! One-window glass just-now break by one-man angry.}
{60}{42}{30.0}
{The warning idiom becomes a single imperative, and the passive stays compact.}

\corpusExample{10}
{She must give her book to him before she can start working.}
{s1must giv s's buc tu h bifor s1can start tu uorc.}
{She-now-must give her book to he before she-now-can start [Object] one-work.}
{59}{50}{15.3}
{Modal structure stays regular, but the sentence keeps more semantic content than the highest-gain cases.}

\subsection{Conditionals, Logic \& Commands}

\corpusExample{11}
{If you don't stop running, you will fall and break your legs.}
{if u1no stop tu run, u1d fol an breic tu u's 2leg.}
{If you-now-not stop [Obj] run, you-in-one-day fall and break [Obj] your two-leg.}
{61}{50}{18.0}
{Conditionals remain explicit, but the clause still carries a lot of event content.}

\corpusExample{12}
{Every morning I drink two cups of coffee and eat one piece of toast before work.}
{*morning i*drinc 2cup cofi an *iit piis toust bifor uorc.}
{Any-morning I-habitually-drink two-cup coffee and habitually-eat one-piece toast before work.}
{80}{57}{28.7}
{Habitual markers reduce repeated aspect information across coordinated verbs.}

\corpusExample{13}
{Take three eggs, add two cups of flour, and mix everything together for ten minutes.}
{teic 3eg, ad 2cup flaur, an mics *ting for 10m.}
{Take three-egg, add two-cup flour, and mix all-thing for ten-minutes.}
{84}{47}{44.0}
{Imperatives and generalized quantity tokens compress recipe-style instructions well.}

\corpusExample{14}
{Please give me two glasses of water, I am very thirsty.}
{~u giv tu i 2glas uoter, i bi >tersti.}
{Polite-you-give [Obj] I two-glass water, I be very-thirsty.}
{55}{38}{30.9}
{A deferential pronoun plus intensity marker replaces polite English padding.}

\subsection{Dialogue \& Description}

\corpusExample{15}
{\"How many children do you have?\" \"I have two boys and one girl.\"}
{"haum ciald u hav?" "i hav 2boi an gerl."}
{"How-many child you have?" "I have two-boy and one-girl."}
{64}{41}{35.9}
{Question words and bare number-nouns remove articles and plural marking cleanly.}

\corpusExample{16}
{Last year, she traveled to five different countries and learned three new languages because she wanted to understand the world better.}
{s-1Y travel tu 5cuntri diferent an lern 3languij niu cos s-1Y uant understand uerld mor gud.}
{She one-year-ago travel to five-country different and learn three-language new because she one-year-ago want understand one-world more good.}
{134}{92}{31.3}
{Long narrative content still benefits from temporal packing and reduced function words.}

\corpusExample{17}
{I might go to the party, but she loves him and he does not love her.}
{i1d ~can gou tu parti, but s1luv h an h no luv s.}
{I in-future might go to one-party, but she-now-love he and he not love she.}
{68}{49}{27.9}
{Probabilistic modality and regular negation compress a multi-clause sentence into a short linear form.}

\corpusExample{18}
{My brother works at a hospital and my sister studies at the university.}
{i's bruder *uorc @ hospital an i's sister *studi @ iuniversiti.}
{My brother any-work at one-hospital and my sister any-study at one-university.}
{71}{63}{11.3}
{Compression stays low when the sentence depends on longer academic or Latinate roots.}

\corpusExample{19}
{The children were playing in the park when it started raining heavily.}
{*ciald plei in parc uen t-0s start rein >hevi.}
{Any-child play in one-park when it-just-now start rain very-heavily.}
{70}{46}{34.3}
{Preposition-as-verb structure and intensity marking remove multiple English helper words.}

\corpusExample{20}
{How much does this car cost? It is too expensive for me.}
{haum dis car cost? t bi >>ecspensiv for i.}
{How-much this car cost? It be extremely-expensive for I.}
{56}{42}{25.0}
{The syntax is regular, but the sentence still carries a relatively dense lexical load.}

\subsection{Advanced Domains}

\corpusExample{21}
{The spaceship launched into orbit after the countdown reached zero.}
{speisship -0s lonci in orbit after caunt-daun -0s hic 0.}
{One-spaceship just-now launch in one-orbit after one-count-down just-now hit zero.}
{67}{56}{16.4}
{Compound formation and numeric notation compress technical narrative moderately well.}

\corpusExample{22}
{To be or not to be, that is the question we must all answer eventually.}
{bi or no bi, dat bi cuesciun *i must anser in taim.}
{Be or not be, that be one-question all-of-us must answer in future-time.}
{71}{51}{28.2}
{Logical operators and universal grouping reduce philosophical phrasing without altering structure.}

\corpusExample{23}
{Our company needs to optimize its supply chain to maximize profit margins next year.}
{*i's compani niid optimaiz t's suplai-cein tu macsimaiz *marjin-profit 1Y.}
{All-of-us-possessive company need optimize its supply-chain to maximize any-margin-profit next-year.}
{84}{74}{11.9}
{Business vocabulary retains much of its length, so the gain comes mainly from grammar compression.}

\corpusExample{24}
{The president announced a new tax policy that will affect middle class families.}
{prezident -0s anauns tacs-polisi niu dat 1d apect *pamili clas-midl.}
{One-president just-now announce one-tax-policy new that in-future affect any-family class-middle.}
{80}{68}{15.0}
{The sentence stays readable, but domain terms limit the total reduction.}

\corpusExample{25}
{The doctor prescribed antibiotics to treat the patient's severe bacterial infection.}
{doctor -0s prescraib tu antibiotik tu triit peishent's infecshun bacteriel >siirius.}
{One-doctor just-now prescribe [Obj] one-antibiotic to treat one-patient's infection bacterial very-serious.}
{84}{84}{0.0}
{Medical vocabulary is already compact in English, so regularization adds only a small gain.}

\corpusExample{26}
{The three bedroom house with a large backyard is currently on the market for rent.}
{haus 3bedruum uid baciard big in marcet for rent nau.}
{One-house three-bedroom with one-backyard big now-in one-market for rent now.}
{82}{53}{35.4}
{Compound nouns and the preposition-verb pattern remove several English helpers at once.}

\corpusExample{27}
{If the server crashes, the backup system will automatically restore the database.}
{if server craish, backup sistem 1d otoristor daatabeis.}
{If one-server crash, one-backup-system in-future auto-restore one-database.}
{81}{55}{32.1}
{The future operator and compound technical vocabulary shorten a longer procedural sentence.}

\corpusExample{28}
{Global warming is melting the polar ice caps faster than scientists predicted.}
{gloubal uorming melt *cap-ais poular >fast dan *saientist -1d pridiict.}
{Global-warming now-melt any-cap-ice polar more-fast than any-scientist past-predict.}
{78}{71}{9.0}
{Scientific phrasing compresses only slightly because most of the content sits in core roots.}

\corpusExample{29}
{The wind whispered through the ancient trees as the silver moon rose in the sky.}
{uind -0s uisper tru *trii >ould az muun silver -0s raiz in scai.}
{One-wind just-now whisper through any-tree very-old as one-moon silver just-now rise in one-sky.}
{80}{64}{20.0}
{Poetic description still benefits from post-modifiers and compact event framing.}

\corpusExample{30}
{The delivery truck will arrive with your package between noon and two o'clock.}
{diliveri truc 1d ariv uid u's pacij btuiin 12:00 an 14:00.}
{One-delivery-truck in-future arrive with your package between 12:00 and 14:00.}
{78}{58}{25.6}
{Range notation and future marking reduce a long scheduling phrase into short tokens.}

\corpusExample{31}
{The defendant formally pleaded not guilty to all charges presented by the court.}
{diipendant -0s pliid formal no gilti tu *carj prizent bai cort.}
{One-defendant just-now plead formally not-guilty to any-charge presented by one-court.}
{80}{63}{21.2}
{Legal transliterations remain relatively long, so most savings come from grammar.}

\corpusExample{32}
{Boil the water, add a pinch of salt, and stir constantly until the sauce thickens.}
{boil tu uoter, ad tu pinci solt, an *ster until sos >tic.}
{Boil [Obj] one-water, add [Obj] one-pinch salt, and any-stir until one-sauce very-thick.}
{82}{57}{30.5}
{Imperatives and generic aspect markers compress procedural language effectively.}

\corpusExample{33}
{I have never felt this way about anyone before, you changed my entire life.}
{i no *fiil dis ue abaut *person bifor, u -0s ceinj i's laif >hol.}
{I not habitually-feel this way about any-person before, you just-now change my life very-whole.}
{75}{65}{13.3}
{Emotive language gains less because it relies on many core semantic roots.}

\corpusExample{34}
{Go straight for two miles, then turn left at the second traffic light you see.}
{gou streit for 2mail, den tern left @ trapic lait tuu u1sii.}
{Go straight for two-mile, then turn left at one-traffic-light two you-now-see.}
{78}{60}{23.1}
{Direction commands still compress through number notation and reduced function words.}

\corpusExample{35}
{Hyperinflation destroys purchasing power and severely halts economic development.}
{haiperinflashun *distroi tu pauer-bai an >siirius holt tu divelopment economik.}
{Hyperinflation habitually-destroy [Obj] power-buy and very-serious halt [Obj] development economic.}
{81}{79}{2.5}
{Low compression reflects the need to preserve long international academic vocabulary.}

%---
\section{Information Density Analysis}

\begin{table}[h!]
\centering
\MEtablelayout
\begin{tabularx}{\linewidth}{lrX}
\rowcolor{MEGray}
\textbf{Example} & \textbf{Reduction} & \textbf{Key Mechanism} \\
Ex.~6: Timestamp     & 67.7\% & ISO 8601 replaces long alphabetic date phrasing \\
Ex.~5: Approximation & 58.6\% & \me{\~{}50person} vs.\ ``About fifty people'' \\
Ex.~4/8: Future packing & 52--53\% & tense prefixing plus compact pronoun/time notation \\
\midrule
Ex.~25: Medical      &  0.0\% & technical vocabulary erases most compression gains \\
Ex.~35: Macroecon.   &  2.5\% & long international academic roots dominate the line \\
\bottomrule
\end{tabularx}
\caption{Compression Extremes (n=35 corpus)}
\end{table}

\textbf{Corpus Density Metric:}
\begin{itemize}[leftmargin=*,noitemsep]
    \item Total English characters: \textbf{2,546}
    \item Total MinEnglish characters: \textbf{1,827}
    \item \textbf{Average reduction: 28.2\%}
\end{itemize}

The maximum gains occur where English relies on auxiliary verb chains and
prepositional phrases. Minimum gains occur with Latinate terminology that
resists phonetic compression.

%---
\section{Theoretical Limitations: The Human Biological Factor}

The structural additions of V1.0.0 (Literalization \me{\_}, Clause Anchors
\me{co}/\me{oc}, Echo-Tagging) address standard linguistic inefficiencies.
However, the primary limitation of MinEnglish is not structural but biological.

Mammalian communicative pathways evolved to facilitate social cohesion,
hierarchical navigation, and emotional subtext through productive
ambiguity$^{[2]}$. MinEnglish enforces an objective, algebraic paradigm.
The absolute requirement for temporal precision, numerical exactness, and
the prohibition of metaphorical obfuscation inherently conflicts with natural
human cognitive behavior$^{[3]}$. Widespread adoption would likely prompt
measurable psychological resistance, as the protocol disables the linguistic
mechanisms humans rely on to manage sociological complexity.

%---
\section{Conclusion}

MinEnglish V1.0.0 establishes a stable baseline for the project: a deliberately
regularized English-derived system built around explicit quantity, explicit
time, phonetic simplification, and operator-based syntax. Across the 35-sentence
comparative corpus, the specification demonstrates that these design choices can
produce substantial character-count reduction while preserving a recognizable
relationship to standard English.

Just as importantly, the current release clarifies the project's real position.
MinEnglish is best understood not as a finished universal language replacement,
but as a formalized communication protocol and compression-oriented grammar
experiment. Its strongest value lies in showing how much of English can be made
regular, explicit, and computationally tractable without discarding the core
lexical shell.

V1.0.0 therefore serves as both a usable specification and a foundation for the
next stage of refinement. Future versions should focus on validating the system
more broadly, reducing remaining ambiguities, and sharpening the boundary between
the minimal core grammar and its optional compression extensions.

%---
\section*{References}
\begin{enumerate}
\item Shannon, C. E. (1948). A mathematical theory of communication. \textit{Bell System Technical Journal}, 27(3), 379--423.
\item Miller, G. A. (1956). The magical number seven, plus or minus two. \textit{Psychological Review}, 63(2), 81--97.
\item Chomsky, N. (1957). \textit{Syntactic Structures}. Mouton.
\item Zipf, G. K. (1949). \textit{Human Behavior and the Principle of Least Effort}. Addison-Wesley.
\end{enumerate}

%---
\begin{appendices}

\section*{MIT License}

\begin{lstlisting}
Copyright (c) 2026 Dan Micsa, PhD

Permission is hereby granted, free of charge, to any person
obtaining a copy of this language specification and associated
documentation files (the "Specification"), to deal in the
Specification without restriction, including without limitation
the rights to use, copy, modify, merge, publish, distribute,
sublicense, and/or sell copies of the Specification, and to
permit persons to whom the Specification is furnished to do so,
subject to the following conditions:

The above copyright notice and this permission notice shall be
included in all copies or substantial portions of the Specification.

THE SPECIFICATION IS PROVIDED "AS IS", WITHOUT WARRANTY OF ANY
KIND. IN NO EVENT SHALL THE AUTHORS BE LIABLE FOR ANY CLAIM,
DAMAGES OR OTHER LIABILITY HOWEVER CAUSED.
\end{lstlisting}

\end{appendices}

\end{document}
